\chapter{Medidas e perímetro}\label{ch:g4:measure}

O Ano 3 mediu com réguas, balanças e relógios
(\cref{ch:g3:measure}); este capítulo acrescenta o quadro de unidades
--- a máquina que converte tudo --- e duas maneiras de medir uma
forma plana: o comprimento do contorno (\omterm{def:g4:measure:perimeter}{perímetro}) e o espaço de
dentro (\omterm{def:g4:measure:area}{área}, contada em quadradinhos por enquanto).

\section{O quadro de unidades}

\begin{method}[Converter com o quadro]\label{met:g4:measure:table}
Para os comprimentos, cada coluna vale dez da seguinte:
\begin{center}
\renewcommand{\arraystretch}{1.3}
\begin{tabular}{c|c|c|c|c|c|c}
km & hm & dam & m & dm & cm & mm \\
\hline
 & & & $3$ & $4$ & $0$ & \\
\end{tabular}
\end{center}
\begin{enumerate}
  \item Escreva o número com um algarismo por coluna, pondo o
        algarismo das unidades da medida na coluna da unidade dela
        (aqui $3.4$ m: o $3$ embaixo de m);
  \item leia o número na nova unidade: $3.4$ m $= 340$ cm (preencha
        com \omterm{def:g1:counting:zero}{zeros} as colunas vazias);
  \item as \omterm{def:g2:measure:units}{massas} (kg\,\dots\,g) e as \omterm{def:g2:measure:units}{capacidades} (L\,\dots\,cL)
        funcionam do mesmo jeito.
\end{enumerate}
\end{method}

\begin{example}\label{ex:g4:measure:convert}
$5$ km $= 5\,000$ m; \quad $270$ cm $= 2.7$ m (o $2$ cai embaixo de
m); \quad $3$ kg $250$ g $= 3\,250$ g; \quad
$2.5$ L $= 250$ cL.
\end{example}

\section{Perímetro}

\begin{definition}[Perímetro]\label{def:g4:measure:perimeter}
O \emph{perímetro}\index{perímetro} de um \omterm{def:g4:geometry:polygon}{polígono} é o comprimento
total do contorno dele: some os comprimentos de todos os seus lados
(na mesma unidade!). Atalhos para os clássicos:
\[
\text{retângulo: } P = 2 \times (L + w),
\qquad
\text{quadrado: } P = 4 \times c ,
\]
em que $L$ é o comprimento, $w$ a largura e $c$ o lado.
\end{definition}

\begin{example}\label{ex:g4:measure:perimeter}
Um jardim retangular de $12$ m por $7$ m:
\[
P = 2 \times (12 + 7) = 2 \times 19 = 38 \text{ m}
\]
de cerca. Uma moldura quadrada de lado $18$ cm:
$P = 4 \times 18 = 72$ cm.
\end{example}

\section{Área: contar quadradinhos}

\begin{definition}[Área por contagem]\label{def:g4:measure:area}
A \emph{área}\index{área} de uma figura desenhada numa malha é o
número de quadradinhos que ela cobre. Dois meios quadradinhos contam
como um. A área responde ``quanta superfície''; o \omterm{def:g4:measure:perimeter}{perímetro} responde
``que comprimento tem o contorno'' --- duas perguntas diferentes!
\end{definition}

\begin{omfigure}
\begin{tikzpicture}[scale=0.55]
  \draw[gray!40, thin] (0,0) grid (13,4);
  \draw[very thick, omDef, fill=omDef!15] (1,1) rectangle (5,3);
  \node[font=\small, omDef] at (3,2) {$8$ quadradinhos};
  \draw[very thick, omThm, fill=omThm!15]
    (7,1) -- (11,1) -- (11,2) -- (9,2) -- (9,3) -- (7,3) -- cycle;
  \node[font=\small, omThm] at (8.4,1.5) {$6$ quadradinhos};
\end{tikzpicture}

{\small Duas figuras na malha: \omterm{def:g4:measure:area}{áreas} de $8$ e $6$ quadradinhos. Agora
percorra os contornos e conte as unidades: $12$ nas duas! Mesmo
\omterm{def:g4:measure:perimeter}{perímetro}, \omterm{def:g4:measure:area}{áreas} diferentes.}
\end{omfigure}

\begin{example}[Meios quadradinhos]\label{ex:g4:measure:half}
Um triângulo retângulo desenhado na malha com catetos $4$ e $2$ cobre
$3$ quadradinhos inteiros e $2$ meios quadradinhos\,\dots\ conte com
cuidado: a \omterm{def:g4:measure:area}{área} dele é a \omterm{def:g3:division:half}{metade} do retângulo $4 \times 2$, ou seja,
$4$ quadradinhos. Cortar um retângulo pela diagonal sempre dá dois
triângulos de mesma \omterm{def:g4:measure:area}{área}.
\end{example}

\begin{example}[Mesma área, perímetros diferentes]\label{ex:g4:measure:samearea}
Um quadrado $4 \times 4$ e um retângulo $8 \times 2$ cobrem os mesmos
$16$ quadradinhos --- mesma \omterm{def:g4:measure:area}{área}. Contornos: $16$ unidades para o
quadrado e $20$ para o retângulo. Nenhuma das duas medidas determina
a outra: uma forma comprida e fina tem muito contorno para pouca
superfície.
\end{example}

\section{Durações}

\begin{method}[Somar horários atravessando a hora]\label{met:g4:measure:time}
Para somar ou subtrair durações, trabalhe em pulos passando pelas
horas cheias (lembre que $1$ h $= 60$ min):
\[
9\text{ h }40 + 35 \text{ min}:
\quad 9\text{ h }40 \xrightarrow{+20} 10\text{ h }00
\xrightarrow{+15} 10\text{ h }15 .
\]
Nunca some minutos passando de $60$ sem converter: $40 + 35 = 75$ min
$= 1$ h $15$ min.
\end{method}

\begin{example}\label{ex:g4:measure:timetable}
O trem parte às $8{:}47$ e chega às $11{:}05$. Duração: de
$8{:}47 \to 9{:}00$ são $13$ min; de $9{:}00 \to 11{:}00$ são $2$ h;
de $11{:}00 \to 11{:}05$ são $5$ min. Total: $2$ h $18$ min.
\end{example}

\section{Exercícios}

\begin{exercise}[$\star$]\label{exo:g4:measure:1}
Converta com o quadro de unidades: $7$ m em cm; $3.2$ km em m; $85$
mm em cm; $4\,500$ m em km.
\end{exercise}

\begin{exercise}[$\star$]\label{exo:g4:measure:2}
Converta: $2$ kg em g; $1\,800$ g em kg e g; $3.5$ L em cL;
$25$ cL em L.
\end{exercise}

\begin{exercise}[$\star$]\label{exo:g4:measure:3}
Calcule o \omterm{def:g4:measure:perimeter}{perímetro}: um retângulo $9$ cm $\times$ $5$ cm; um quadrado
de lado $7.5$ cm; um triângulo de lados $6$ cm, $8$ cm e $10$ cm.
\end{exercise}

\begin{exercise}[$\star$]\label{exo:g4:measure:4}
Um retângulo tem \omterm{def:g4:measure:perimeter}{perímetro} $26$ cm e comprimento $8$ cm. Encontre a
largura dele (ache primeiro comprimento $+$ largura).
\end{exercise}

\begin{exercise}[$\star$]\label{exo:g4:measure:5}
No papel quadriculado, desenhe uma figura em forma de L que cubra
exatamente $10$ quadradinhos e conte as unidades do \omterm{def:g4:measure:perimeter}{perímetro} dela.
\end{exercise}

\begin{exercise}[$\star$]\label{exo:g4:measure:6}
Trace dois retângulos diferentes de \omterm{def:g4:measure:area}{área} $12$ quadradinhos. Calcule o
\omterm{def:g4:measure:perimeter}{perímetro} de cada um. Qual é o mais ``redondo'' e qual o mais
``fininho''?
\end{exercise}

\begin{exercise}[$\star$]\label{exo:g4:measure:7}
Um triângulo retângulo é desenhado na malha com catetos $6$ e $3$.
Qual é a \omterm{def:g4:measure:area}{área} dele em quadradinhos? (Use o
\cref{ex:g4:measure:half}.)
\end{exercise}

\begin{exercise}[$\star$]\label{exo:g4:measure:8}
Calcule as durações: de $10{:}20$ até $12{:}45$; de $7{:}35$ até
$9{:}10$. Um filme dura $1$ h $50$ e começa às $20{:}30$: a que
horas ele termina?
\end{exercise}

\begin{exercise}[$\star$]\label{exo:g4:measure:9}
Kim corre $3$ voltas em torno de um campo retangular de $60$ m por
$45$ m. Quantos metros ela corre?
\end{exercise}

\begin{exercise}[$\star\star$]\label{exo:g4:measure:10}
Um agricultor quer cercar um terreno quadrado de lado $85$ m,
deixando um portão de $4$ m sem cerca. Quantos metros de cerca ele
tem de comprar?
\end{exercise}

\begin{exercise}[$\star\star$]\label{exo:g4:measure:11}
Verdadeiro ou falso, com exemplos na malha: ``se duas figuras têm o
mesmo \omterm{def:g4:measure:perimeter}{perímetro}, elas têm a mesma \omterm{def:g4:measure:area}{área}''; ``se duas figuras têm a
mesma \omterm{def:g4:measure:area}{área}, elas têm o mesmo \omterm{def:g4:measure:perimeter}{perímetro}''.
(O \cref{ex:g4:measure:samearea} e a figura do capítulo vão
ajudar.)
\end{exercise}
